\subsection{UNIQUENESS THEOREM}

PROPOSITION 4. Let \((\mathfrak{g},\mathfrak{h},\mathrm{B},(X_{\alpha})_{\alpha\in\mathrm{B}})\) be a framed semi-simple Lie algebra. Let \((n(\alpha,\beta))_{\alpha,\beta\in\mathrm{B}}\) and \((X_{\alpha})_{\alpha\in\mathrm{B}\cup(-\mathrm{B})}\) be the corresponding Cartan matrix and generating family.

\begin{enumerate}
\def\labelenumi{(\roman{enumi})}
\item
  The family \(((X_{\alpha}, H_{\alpha}, X_{-\alpha}))_{\alpha \in \mathrm{B}}\) and the relations (16) to (22) of no. 3 constitute a presentation of \(\mathfrak{g}\) .
\item
  The family \((X_{\alpha})_{\alpha \in \mathrm{B}}\) and the relations (21) of no. 3 constitute a presentation of the subalgebra of \(\mathfrak{g}\) generated by \((X_{\alpha})_{\alpha \in \mathrm{B}}\) .
\end{enumerate}

Let R be the root system of \((\mathfrak{g},\mathfrak{h})\) . Applying to R and B the constructions of nos. 2 and 3, we obtain objects that we shall denote by \(a', g', X'_{\alpha}, H'_{\alpha}, \ldots\) instead of \(a, g, X_{\alpha}, H_{\alpha}, \ldots\)

There exists a homomorphism \(\varphi\) from the Lie algebra \(\mathfrak{g}'\) to the Lie algebra \(\mathfrak{g}\) such that \(\varphi(X_{\alpha}') = X_{\alpha}\) , \(\varphi(H_{\alpha}') = H_{\alpha}\) , \(\varphi(X_{-\alpha}') = X_{-\alpha}\) for all \(\alpha \in B\) (Prop. 1). Since \(\dim \mathfrak{g}' = \operatorname{Card} R + \operatorname{Card} B = \dim \mathfrak{g}\) , \(\varphi\) is bijective. This proves (i).

The subalgebra of \(\mathfrak{g}' = \mathfrak{a}' / (\mathfrak{n}'\oplus \theta'\mathfrak{n}') = (\mathfrak{a}_+^\prime \oplus \mathfrak{a}^{\prime 0}\oplus \mathfrak{a}_-^\prime) / (\mathfrak{n}'\oplus \theta'\mathfrak{n}')\) generated by \((X_{\alpha}^{\prime})_{\alpha \in \mathrm{B}}\) can be identified with \(\mathfrak{a}_+^\prime /\mathfrak{n}'\) . In view of Prop. 3 and the definition of \(\mathfrak{n}'\) , this proves (ii).

COROLLARY. Every framed semi-simple Lie algebra is obtained from a framed semi-simple \(\mathbf{Q}\) -Lie algebra by extension of scalars from \(\mathbf{Q}\) to \(k\) .

This follows immediately from the proposition.

THEOREM 2. Let \((\mathfrak{g},\mathfrak{h},\mathrm{B},(X_{\alpha})_{\alpha \in \mathrm{B}})\) and \((\mathfrak{g}',\mathfrak{h}',\mathrm{B}',(X'_{\alpha})_{\alpha \in \mathrm{B}'})\) be framed semi-simple Lie algebras, let R and R' be the root systems of \((\mathfrak{g},\mathfrak{h})\) and \((\mathfrak{g}',\mathfrak{h}')\) , let \((n(\alpha ,\beta))_{\alpha ,\beta \in \mathrm{B}}\) (resp. \((n^{\prime}(\alpha ,\beta))_{\alpha ,\beta \in \mathrm{B}'}\) ) be the Cartan matrix of R (resp. R') relative to B (resp. B'), and let \(\Delta\) (resp. \(\Delta'\) ) be the Dynkin graph of R (resp. R') relative to B (resp. B').

\begin{enumerate}
\def\labelenumi{(\roman{enumi})}
\item
  If \(\varphi\) is an isomorphism from \(\mathfrak{h}^*\) to \(\mathfrak{h}'^*\) such that \(\varphi(\mathrm{R}) = \mathrm{R}'\) and \(\varphi(\mathrm{B}) = \mathrm{B}'\) , there exists a unique isomorphism \(\psi\) from \((\mathfrak{g}, \mathfrak{h}, \mathrm{B}, (X_{\alpha})_{\alpha \in \mathrm{B}})\) to \((\mathfrak{g}', \mathfrak{h}', \mathrm{B}', (X_{\alpha}')_{\alpha \in \mathrm{B}'})\) such that \(\psi|_{\mathfrak{h}} = {}^t\varphi^{-1}\) .
\item
  If \(f\) is a bijection from \(\mathrm{B}\) to \(\mathrm{B}'\) such that \(n'(f(\alpha), f(\beta)) = n(\alpha, \beta)\) for all \(\alpha, \beta \in \mathrm{B}\) , there exists an isomorphism from \((\mathfrak{g}, \mathfrak{h}, \mathrm{B}, (X_{\alpha})_{\alpha \in \mathrm{B}})\) to \((\mathfrak{g}', \mathfrak{h}', \mathrm{B}', (X_{\alpha}')_{\alpha \in \mathrm{B}'}).\)
\item
  If there exists an isomorphism from \(\Delta\) to \(\Delta'\) , there exists an isomorphism from \((\mathfrak{g},\mathfrak{h},\mathrm{B},(X_{\alpha})_{\alpha \in \mathrm{B}})\) to \((\mathfrak{g}',\mathfrak{h}',\mathrm{B}',(X_{\alpha}')_{\alpha \in \mathrm{B}'})\) .
\end{enumerate}

This follows immediately from Prop. 4 (i) (making use of Chap. VI, §4, no. 2, Prop. 1 for part (iii)).

Scholium. To any splittable semi-simple Lie algebra g is associated a Dynkin graph, which determines g up to isomorphism (Th. 2 (iii)). This graph is non-empty and connected if and only if g is simple (§3, no. 2, Cor. 1 of Prop. 6). By Th. 1 of no. 3, and Chap. VI, §4, no. 2, Th. 3, the splittable simple Lie algebras are the algebras of type \(A_{l}\) ( \(l \geq 1\) ), \(B_{l}\) ( \(l \geq 2\) ), \(C_{l}\) ( \(l \geq 3\) ), \(D_{l}\) ( \(l \geq 4\) ), \(E_{6}\) , \(E_{7}\) , \(E_{8}\) , \(F_{4}\) , \(G_{2}\) . No two algebras in this list are isomorphic.

PROPOSITION 5. Let \((\mathfrak{g},\mathfrak{h},\mathrm{B},(X_{\alpha})_{\alpha \in \mathrm{B}})\) be a framed semi-simple Lie algebra, and \((X_{\alpha})_{\alpha \in \mathrm{B}\cup (-\mathrm{B})}\) the corresponding generating family. There exists a unique automorphism \(\theta\) of \(\mathfrak{g}\) such that \(\theta (X_{\alpha}) = X_{-\alpha}\) for all \(\alpha \in \mathrm{B}\cup (-\mathrm{B})\) . We have \(\theta^2 = \mathrm{Id}_{\mathfrak{g}}\) , and \(\theta (h) = -h\) for all \(h\in \mathfrak{h}\) .

The uniqueness is clear since \((X_{\alpha})_{\alpha\in\mathrm{B}\cup(-\mathrm{B})}\) generates the Lie algebra g. In view of Prop. 4, the existence of \(\theta\) follows from what we said in no. 3 before Th. 1.

COROLLARY. Let \((\mathfrak{g},\mathfrak{h})\) be a split semi-simple Lie algebra. Then \((\mathfrak{g},\mathfrak{h})\) possesses a Chevalley system (§2, no. 4, Def. 3).

Let R be the root system of \((\mathfrak{g},\mathfrak{h})\) . For all \(\alpha\in R\) , let \(X_{\alpha}\) be a non-zero element of \(g^{\alpha}\) . Assume that the \(X_{\alpha}\) are chosen so that \([X_{\alpha},X_{-\alpha}]=-H_{\alpha}\) for all \(\alpha\in R\) (§2, no. 4, Lemma 2). Let B be a basis of R and \(\theta\) the automorphism of g such that \(\theta(X_{\alpha})=X_{-\alpha}\) for all \(\alpha\in\mathrm{B}\cup(-\mathrm{B})\) . We have \(\theta|_{\mathfrak{h}}=-\mathrm{Id}_{\mathfrak{h}}\) . Hence, for all \(\alpha\in R\) there exists \(t_{\alpha}\in k^{*}\) such that \(\theta X_{\alpha}=t_{\alpha}X_{-\alpha}\) . We have

\[
\begin{array}{c} t _ {\alpha} t _ {- \alpha} H _ {\alpha} = [ t _ {\alpha} X _ {- \alpha}, t _ {- \alpha} X _ {\alpha} ] = [ \theta X _ {\alpha}, \theta X _ {- \alpha} ] = \theta ([ X _ {\alpha}, X _ {- \alpha} ]) \\ = \theta (- H _ {\alpha}) = H _ {\alpha} \end{array}
\]

so \(t_{\alpha}t_{-\alpha} = 1\) for all \(\alpha \in \mathrm{R}\) . Introduce the \(\mathrm{N}_{\alpha \beta}\) as in §2, no. 4. If \(\alpha, \beta, \alpha + \beta \in \mathrm{R}\) ,

\[
\begin{array}{r} \mathrm{N} _ {- \alpha , - \beta} t _ {\alpha} t _ {\beta} X _ {- \alpha - \beta} = t _ {\alpha} t _ {\beta} [ X _ {- \alpha}, X _ {- \beta} ] = [ \theta X _ {\alpha}, \theta X _ {\beta} ] = \theta ([ X _ {\alpha}, X _ {\beta} ]) \\ = \mathrm{N} _ {\alpha \beta} \theta X _ {\alpha + \beta} = \mathrm{N} _ {\alpha \beta} t _ {\alpha + \beta} X _ {- \alpha - \beta} \end{array}
\]

so, in view of §2, no. 4, Lemma 4,

\[
(q + 1) ^ {2} t _ {\alpha} t _ {\beta} = \mathrm{N} _ {\alpha \beta} ^ {2} t _ {\alpha + \beta}
\]

where \(q\) is an integer. It follows that if \(t_{\alpha}\) and \(t_{\beta}\) are squares in \(k^*\) , so is \(t_{\alpha + \beta}\) . Since \(t_{\alpha} = 1\) for all \(\alpha \in \mathrm{B}\) , Prop. 19 of Chap. VI, §1, no. 6, proves that \(t_{\alpha}\) is a square for all \(\alpha \in \mathrm{R}\) . Choose, for all \(\alpha \in \mathrm{R}\) , a \(u_{\alpha} \in k\) such that \(u_{\alpha}^{2} = t_{\alpha}\) . This choice can be made so that \(u_{\alpha}u_{-\alpha} = 1\) for all \(\alpha \in \mathrm{R}\) . Put \(X_{\alpha}' = u_{\alpha}^{-1}X_{\alpha}\) . Then, for all \(\alpha \in \mathrm{R}\) ,

\[
X _ {\alpha} ^ {\prime} \in \mathfrak {g} ^ {\alpha}, [ X _ {\alpha} ^ {\prime}, X _ {- \alpha} ^ {\prime} ] = [ X _ {\alpha}, X _ {- \alpha} ] = - H _ {\alpha},
\]

and \(\theta(X_{\alpha}^{\prime}) = \theta(u_{\alpha}^{-1}X_{\alpha}) = u_{\alpha}^{-1}t_{\alpha}X_{-\alpha} = u_{\alpha}X_{-\alpha} = u_{\alpha}u_{-\alpha}X_{-\alpha}^{\prime} = X_{-\alpha}^{\prime}\) , so that \((X_{\alpha}^{\prime})_{\alpha \in \mathrm{R}}\) is a Chevalley system of \((\mathfrak{g},\mathfrak{h})\) .

